\section{Lexical probabilities}
We assume a lexical probability assignment
\[
\Lex : \Sigma \times \At \to [0,1]
\]
such that for every word $w \in \Sigma$,
\[
\sum_{p \in \At} \Lex(w,p) = 1,
\]
where all but finitely many summands are zero. Intuitively, $\Lex(w,p)$ is the probability that
the word $w$ is assigned the atomic type $p$.

\begin{remark}
If a word $w \in \Sigma$ occurs in the corpus but is never assigned an atomic type $p$, then
$\Lex(w,p)=0$. This matches the intuition from the draft: lack of attestation for a type within
the observed vocabulary should count against that type.
\end{remark}

\begin{remark}
Words outside the corpus vocabulary are not assigned probabilities here at all. Rather than
assigning them an arbitrary value such as $0$ or leaving them as exceptional elements with
undefined behaviour, they are simply excluded from the domain of the semantics. This keeps the
formal system clean and makes explicit that out-of-vocabulary handling is a future extension.
\end{remark}

